% Generated draft: requires mathematical and source review.
\documentclass[11pt]{article}
\usepackage[margin=25mm]{geometry}
\usepackage{fontspec}
\setmainfont{Latin Modern Roman}
\usepackage{amsmath,amssymb,mathtools,unicode-math}
\setmathfont{Latin Modern Math}
\usepackage{enumitem,needspace,xurl,hyperref}
\setlength{\parindent}{0pt}
\setlength{\parskip}{5pt}
\setlength{\emergencystretch}{3em}
\newcommand{\N}{\mathbb{N}}
\newcommand{\Z}{\mathbb{Z}}
\newcommand{\Q}{\mathbb{Q}}
\newcommand{\R}{\mathbb{R}}
\newcommand{\C}{\mathbb{C}}
\newcommand{\F}{\mathbb{F}}
\newcommand{\A}{\mathbb{A}}
\newcommand{\PP}{\mathbb{P}}
\newcommand{\E}{\mathbb{E}}
\newcommand{\eps}{\varepsilon}
\newcommand{\dd}{\,\mathrm{d}}
\begin{document}
\section*{UnsolvedMath ID: 4000014}
\textbf{Source problem code:} catalogue entry
\subsection{Definitions and mathematical statement}
Let $X$ be a state space, $d_0$ an initial metric, and $m_x^t$ the transition probability from $x$ at time $t$. For a metric $d_t$ for which these measures have finite first moments, define
\[
W_{1,d_t}(\mu,\nu)=\inf_{\pi\in\Pi(\mu,\nu)}\int d_t(u,v)\,d\pi(u,v),\qquad
\kappa_t(x,y)=1-\frac{W_{1,d_t}(m_x^t,m_y^t)}{d_t(x,y)}\quad(x\ne y).
\]
The proposed discrete Ricci flow is the pairwise evolution
\[
\frac{d}{dt}d_t(x,y)=-\kappa_t(x,y)d_t(x,y)
=W_{1,d_t}(m_x^t,m_y^t)-d_t(x,y),\qquad d_t|_{t=0}=d_0.
\]
One branch keeps the kernel fixed, $m_x^t=m_x$. Another permits a suitably evolving kernel coupled to the metric, but the source supplies no evolution law for that kernel.

Determine conditions for existence of this metric-valued flow, its behavior, and its limiting geometry. This includes specifying when the pairwise evolution preserves a nondegenerate metric and determining what long-time conclusions hold in the intended class of Markov spaces. Any normalized or rescaled limiting geometry requires an explicit normalization, which the source does not prescribe.

The supplied summary reports existence, uniqueness, and convergence to constant-curvature metrics for finite-graph versions, while leaving general infinite graphs and jointly evolving kernels unresolved. It refers to ``the above works'' without including those references or their precise flow conventions. Thus this report does not identify which hypotheses, normalizations, or finite-graph variants apply to the displayed pairwise equation.
\subsection{Short English statement}
Evolve distances at a rate given by their current transport-based coarse Ricci curvature. Determine existence, behavior, and limiting geometry, both for a fixed random walk and for a specified coupled evolution of the walk.
\subsection{Sources}
\begin{itemize}
\item Yann Ollivier, \emph{Discrete Ricci curvature: Open problems} (2008), Problem N, \url{https://www.yann-ollivier.org/rech/publs/problems_curvmarkov.pdf}; accessed 2026-07-29, with PDF-text extraction recorded in the supplied provenance.
\item The research summary refers to finite-graph results by ``the above works,'' but no authors, titles, or other identifiers for those works are present in this record.
\item Local registry \texttt{data/unsolvedmath/problems.json}, permanent UnsolvedMath ID $4000014$, source code catalogue entry, \emph{Discrete Ricci flow}. Dataset attribution: Ulam AI and the UnsolvedMath Contributors.
\end{itemize}
\end{document}
